\section*{Hoofdstuk \ref{ch:b2:linalg} --- Lineaire algebra}

\begin{solution}{exo:b2:linalg:1}
Drie vormen in een driedimensionale \omterm{def:b2:linalg:dual}{duale ruimte}: vrijheid
volstaat. Een relatie $\alpha\varphi_1 + \beta\varphi_2 +
\gamma\varphi_3 = 0$ geëvalueerd in $(1,0,0), (0,1,0), (0,0,1)$
geeft $\alpha + \gamma = 0$, $\alpha + \beta = 0$, $\beta + \gamma
= 0$, waaruit $\alpha = \beta = \gamma = 0$.

Antiduale basis $(u_1, u_2, u_3)$: los $\varphi_i(u_j) =
\delta_{ij}$ op. Met $u_j = (x, y, z)$ geeft voor $u_1$ het
stelsel $x + y = 1$, $y + z = 0$, $x + z = 0$ dat $u_1 =
\bigl(\tfrac12, \tfrac12, -\tfrac12\bigr)$; symmetrisch $u_2 =
\bigl(-\tfrac12, \tfrac12, \tfrac12\bigr)$ en $u_3 =
\bigl(\tfrac12, -\tfrac12, \tfrac12\bigr)$.
\end{solution}

\begin{solution}{exo:b2:linalg:2}
Eerste matrix: de enige permutatie met een product ongelijk aan nul
stuurt $1 \mapsto 3$, $2 \mapsto 2$, $3 \mapsto 1$ --- de
\omterm{def:b2:structures:sn}{transpositie} $(1\,3)$, met signatuur $-1$: de \omterm{def:b2:linalg:det}{determinant} is
$-abc$.

Tweede: een permutatie met een product ongelijk aan nul kan de twee
blokken niet mengen (een element dat ze verbindt is $0$), dus
splitst zij in een permutatie van $\{1,2\}$ maal een van
$\{3,4\}$, en de signatuur is het product van beide signaturen: de
som factoriseert als
\[
(ad - bc)(eh - fg) .
\]
De regel die dit suggereert (en die met hetzelfde bewijs juist is):
de \omterm{def:b2:linalg:det}{determinant} van een blokdiagonaalmatrix is het product van de
\omterm{def:b2:linalg:det}{determinanten} van de blokken.
\end{solution}

\begin{solution}{exo:b2:linalg:3}
$F$ wordt bepaald door de twee onafhankelijke vergelijkingen
$\varphi_1(x,y,z,t) = x + y - z - t = 0$ en $\varphi_2 = x - 2y =
0$: leest men \cref{thm:b2:linalg:annihilator} achterstevoren, dan
is $F^\circ = \operatorname{Vect}(\varphi_1, \varphi_2)$ --- zij
liggen per constructie in $F^\circ$, zij zijn vrij (niet
evenredig), en $\dim F^\circ = 4 - \dim F = 4 - 2 = 2$ omdat $\dim
F = 2$ (twee onafhankelijke vergelijkingen in $\R^4$). Basis:
$(\varphi_1, \varphi_2)$; dimensies: $2 + 2 = 4$, zoals de stelling
eist.
\end{solution}

\begin{solution}{exo:b2:linalg:4}
De $\varphi_i$ zijn $n + 1$ vormen op een $(n+1)$-dimensionale
ruimte: vrijheid volstaat. Geldt $\sum_i \lambda_i \varphi_i = 0$,
evalueer dan in de lagrangeveelterm $L_j$ van de knooppunten:
$\lambda_j = 0$. De antiduale basis is $(L_0, \dots, L_n)$, want
$\varphi_i(L_j) = L_j(a_i) = \delta_{ij}$.

Voor de integraalvorm met knooppunten $0, \frac12, 1$ op
$\R_2[X]$: $\int_0^1 P = \sum_i c_i P(a_i)$ met $c_i = \int_0^1
L_i$. Uitgerekend: $L_0 = 2(X - \tfrac12)(X - 1)$, $\int_0^1 L_0 =
\frac16$; $L_1 = -4X(X-1)$, $\int_0^1 L_1 = \frac46$; $L_2 = 2X(X
- \tfrac12)$, $\int_0^1 L_2 = \frac16$. Bijgevolg is
\[
\int_0^1 P = \frac{1}{6}\Bigl(P(0) + 4P\bigl(\tfrac12\bigr) +
P(1)\Bigr)
\quad (P \in \R_2[X]) :
\]
de regel van Simpson, exact op tweedegraadsveeltermen --- een
uitspraak over \omterm{def:b2:linalg:dual}{duale} basissen.
\end{solution}

\begin{solution}{exo:b2:linalg:5}
Er geldt $\operatorname{im} u = Ka$ voor zekere $a \neq 0$; dan is
$u(x) = \varphi(x)\,a$, waarbij $\varphi(x)$ de coördinaat van
$u(x)$ op $a$ is --- lineair in $x$. Spoor: vul $a = e_1$ aan tot
een basis; de matrix van $u$ heeft kolommen $\varphi(e_j)\,e_1$,
dus is haar enige diagonaalelement $\varphi(e_1) = \varphi(a)$:
$\operatorname{tr} u = \varphi(a)$. Vervolgens
\[
u^2(x) = \varphi(x)\, u(a) = \varphi(x)\varphi(a)\, a
= (\operatorname{tr} u)\, u(x).
\]
\omterm{def:b2:linalg:det}{Determinant}, in twee gevallen. \emph{Als $\varphi(a) \neq 0$:} neem
een willekeurige basis van het hypervlak $\ker\varphi$ en voeg $a$
toe. Dan doodt $u$ de ruimte $\ker\varphi$ (daar is $u(x) =
\varphi(x)a = 0$) en is $u(a) = \varphi(a)\,a$: de matrix van $I +
u$ is diagonaal, $(1, \dots, 1,\, 1 + \varphi(a))$, dus $\det(I +
u) = 1 + \varphi(a) = 1 + \operatorname{tr} u$. \emph{Als
$\varphi(a) = 0$:} dan is $a \in \ker\varphi$; neem een basis van
$\ker\varphi$ waarvan de eerste vector $a$ is, en voeg een vector
$b$ met $\varphi(b) = 1$ toe. Dan houdt $I + u$ de basis van
$\ker\varphi$ vast en stuurt zij $b \mapsto b + a$: driehoekig met
diagonaal vol enen, dus $\det(I + u) = 1 = 1 + \operatorname{tr}
u$. Beide gevallen stemmen met de formule overeen.
\end{solution}

\begin{solution}{exo:b2:linalg:6}
Volgens \cref{exo:b2:linalg:9} is het hypervlak $H_A = \{M :
\operatorname{tr}(AM) = 0\}$ met $A \neq 0$.

\emph{Als $A = \lambda I$:} dan is $H_A$ het hypervlak van de
matrices met spoor nul; de matrix van de permutatie die één
$n$-cykel is (enen op de plaatsen $(i, i+1)$ en $(n, 1)$) is
inverteerbaar (haar \omterm{def:b2:linalg:det}{determinant} is $\pm 1$ volgens de berekening
van \cref{ex:b2:linalg:permexample}) en heeft spoor nul.

\emph{Als $A$ niet scalair is:} zoek eerst een inverteerbare $P$
zodanig dat $B = P^{-1}AP$ buiten de diagonaal een element $b_{ji}
\neq 0$ heeft ($j \neq i$). Heeft $A$ er al een, neem dan $P = I$;
is $A$ diagonaal met twee verschillende elementen $d_1 \neq d_2$,
dan levert conjugatie met de transvectie $P = I + E_{12}$ het
element $d_1 - d_2 \neq 0$ buiten de diagonaal op (reken na:
$P^{-1}AP = A + (d_1 - d_2)E_{12}$); en een diagonaalmatrix waarvan
alle elementen gelijk zijn, is scalair en dus uitgesloten. Zet nu
$M' = I + tE_{ij}$ met $t = -\operatorname{tr}(B)/b_{ji}$: dan is
\[
\operatorname{tr}(BM') = \operatorname{tr} B + t\,b_{ji} = 0,
\]
en $M'$ is inverteerbaar (driehoekig met diagonaal vol enen).
Maken we de conjugatie ongedaan, dan is $M = PM'P^{-1}$
inverteerbaar met $\operatorname{tr}(AM) = \operatorname{tr}(BM')
= 0$: dus $M \in H_A$.
\end{solution}

\begin{solution}{exo:b2:linalg:7}
Volgens de permutatieformule is $\det(I + tA)$ een veelterm in $t$;
haar constante term is $1$ (neem $t = 0$). De coëfficiënt van $t$:
ontwikkel $\det$ als een \omterm{def:b2:linalg:alternating}{alternerende vorm} van de kolommen $e_j +
t\,c_j(A)$; wegens de multilineariteit vervangen de termen die
lineair in $t$ zijn precies één $e_j$ door $c_j(A)$:
\[
\sum_{j} \det(e_1, \dots, c_j(A), \dots, e_n)
= \sum_j a_{jj} = \operatorname{tr} A ,
\]
(de \omterm{def:b2:linalg:det}{determinant} met alle canonieke kolommen behalve $c_j(A)$ op
plaats $j$ pikt het $j$-de diagonaalelement op). De afgeleide in
$0$ is dus $\operatorname{tr} A$.

Zij $g(t) = \det(\eu^{tA})$. De groepseigenschap geeft $g(s + t) =
g(s)g(t)$ (de multiplicativiteit van $\det$), $g$ is
differentieerbaar, en $g'(0) = \operatorname{tr} A$ volgens het
voorgaande ($\eu^{tA} = I + tA + O(t^2)$). Een differentieerbaar
morfisme $(\R, +) \to (\R^*, \times)$ voldoet aan $g' = g'(0)\,g$
(differentieer $g(s+t)$ naar $s$ in $0$), dus is $g(t) =
\eu^{t\operatorname{tr} A}$ wegens de eenduidigheid van de
oplossingen van $y' = cy$ met $y(0) = 1$ (volume van bachelorjaar
1).
\end{solution}

\begin{solution}{exo:b2:linalg:8}
Zij $v_k = (1, j^k, j^{2k})^{\mathsf T}$ voor $k = 0, 1, 2$. Met $1
+ j + j^2 = 0$ en $j^3 = 1$:
\[
C v_k =
\begin{pmatrix}
a + b j^k + c j^{2k}\\
c + a j^k + b j^{2k}\\
b + c j^k + a j^{2k}
\end{pmatrix}
= (a + b j^k + c j^{2k})
\begin{pmatrix} 1\\ j^k\\ j^{2k}\end{pmatrix},
\]
(ga de tweede rij na: $j^k(a + bj^k + cj^{2k}) = aj^k + bj^{2k} +
cj^{3k} = c + aj^k + bj^{2k}$). Dus is $v_k$ een eigenvector met
eigenwaarde $\lambda_k = a + bj^k + cj^{2k}$. De $v_k$ vormen een
basis (Vandermonde van de verschillende $1, j, j^2$), dus is $C$
diagonaliseerbaar met deze eigenwaarden en
\[
\det C = \lambda_0\lambda_1\lambda_2
= (a+b+c)(a + bj + cj^2)(a + bj^2 + cj).
\]
\end{solution}

\begin{solution}{exo:b2:linalg:9}
De afbeelding $\Theta \colon A \mapsto \operatorname{tr}(A\,\cdot)$
is lineair van $\mathcal{M}_n(K)$ naar haar \omterm{def:b2:linalg:dual}{duale ruimte}, tussen
ruimten van gelijke dimensie $n^2$: injectiviteit volstaat. Geldt
$\operatorname{tr}(AM) = 0$ voor alle $M$, neem dan $M = E_{ji}$:
$\operatorname{tr}(A E_{ji}) = a_{ij} = 0$ voor alle $i, j$, dus $A
= 0$. Bijgevolg is $\Theta$ een isomorfisme.

Eenduidigheid van het spoor (\cref{prop:b2:linalg:trace}): een vorm
$t$ die alle commutatoren doodt, is $\operatorname{tr}(A\,\cdot)$
voor zekere $A$ met $\operatorname{tr}(A(MN - NM)) = 0$ voor alle
$M, N$, dat wil zeggen $\operatorname{tr}((AM - MA)N) = 0$ voor
alle $N$ (cyclische invariantie), dat wil zeggen $AM = MA$ voor
alle $M$ (injectiviteit van $\Theta$): $A$ commuteert met alles en
is dus scalair ($A$ commuteert met alle $E_{ij}$, wat de elementen
buiten de diagonaal nul maakt en de diagonaalelementen aan elkaar
gelijk), zodat $t = c \operatorname{tr}$.
\end{solution}

\begin{solution}{exo:b2:linalg:10}
Werk over $\C$ (een reële matrix is nilpotent dan en slechts dan
als zij het als complexe matrix is: nilpotentie is $u^n = 0$).

\emph{Stap 1: $\operatorname{tr}(u^k) = 0$ voor $k \geq 1$.} Met
inductie geldt $u^k v - v u^k = k\, u^k$: voor $k = 1$ is dat de
hypothese, en voor de stap
\[
u^{k+1}v - vu^{k+1} = u^k(uv - vu) + (u^k v - v u^k)u
= u^{k+1} + k\,u^{k+1} .
\]
Sporen nemen geeft $0 = \operatorname{tr}(u^k v) -
\operatorname{tr}(vu^k) = k \operatorname{tr}(u^k)$, dus
$\operatorname{tr}(u^k) = 0$.

\emph{Stap 2: sporen van machten nul impliceert nilpotentie (over
$\C$).} Zijn $\lambda_1, \dots, \lambda_r$ de verschillende
eigenwaarden van $u$ ongelijk aan nul, met multipliciteiten $m_1,
\dots, m_r$ (in de karakteristieke veelterm, die over $\C$ in
lineaire factoren uiteenvalt --- \cref{ch:b2:reduction}). De sporen
van de machten zijn $\operatorname{tr}(u^k) = \sum_i m_i
\lambda_i^k$ (trigonaliseer: de diagonaal van de $k$-de macht van
een driehoeksmatrix bestaat uit de $k$-de machten). Het stelsel
$\sum_i m_i \lambda_i^k = 0$ voor $k = 1, \dots, r$ is
Vandermonde-inverteerbaar in de onbekenden $m_i\lambda_i$ (de
matrix $(\lambda_i^{k-1})$ maal de diagonaal $\lambda_i$, met alle
$\lambda_i \neq 0$ verschillend): elke $m_i \lambda_i = 0$, wat
met $m_i \geq 1$ onmogelijk is tenzij $r = 0$. Dus heeft $u$ geen
eigenwaarde ongelijk aan nul: haar karakteristieke veelterm is
$(-X)^n$, en Cayley--Hamilton (\cref{ch:b2:reduction}) geeft $u^n
= 0$: nilpotent.
\end{solution}

\begin{solution}{exo:b2:linalg:11}
Schrijf $V(a_0, \dots, a_n)$ voor de \omterm{def:b2:linalg:det}{determinant} en gebruik
inductie naar $n$; $V(a_0) = 1$ is het beginpunt. Leg $a_0, \dots,
a_{n-1}$ vast en beschouw $D(T) = V(a_0, \dots, a_{n-1}, T)$, de
\omterm{def:b2:linalg:det}{determinant} met laatste kolom $(1, T, \dots, T^n)$: ontwikkeling
langs die kolom laat zien dat $D$ een veelterm van graad $\leq n$
in $T$ is waarvan de coëfficiënt van $T^n$ de minor $V(a_0, \dots,
a_{n-1})$ is. Neem eerst aan dat $a_0, \dots, a_{n-1}$
verschillend zijn. Voor elke $T = a_i$ ($i < n$) vallen twee
kolommen samen, dus $D(a_i) = 0$: met $n$ verschillende nulpunten
en graad $\leq n$ is
\[
D(T) = V(a_0, \dots, a_{n-1}) \prod_{i=0}^{n-1}(T - a_i),
\]
en $T = a_n$ samen met de inductiehypothese geeft de
productformule. Vallen twee van de $a_0, \dots, a_{n-1}$ samen,
dan zijn beide leden $0$ (herhaalde kolommen; een herhaalde
factor) en geldt de formule triviaal.
\end{solution}

\begin{solution}{exo:b2:linalg:12}
\emph{$D$ inverteerbaar.} Vermenigvuldig rechts met de blokmatrix
$T = \begin{psmallmatrix} I & 0\\ -D^{-1}C & I\end{psmallmatrix}$,
die blokdriehoekig is met diagonaal vol eenheidsblokken en $\det T
= 1$ (haar \omterm{def:b2:linalg:det}{determinant} pikt volgens de permutatieformule alleen de
diagonaalblokken op --- de blokregel van \cref{exo:b2:linalg:2}):
\[
\begin{pmatrix} A & B\\ C & D\end{pmatrix} T
= \begin{pmatrix} A - BD^{-1}C & B\\ C - DD^{-1}C & D\end{pmatrix}
= \begin{pmatrix} A - BD^{-1}C & B\\ 0 & D\end{pmatrix},
\]
met \omterm{def:b2:linalg:det}{determinant} $\det(A - BD^{-1}C)\det D = \det\bigl((A -
BD^{-1}C)D\bigr) = \det(AD - BD^{-1}CD)$. Omdat $CD = DC$, is
$BD^{-1}CD = BC$: de \omterm{def:b2:linalg:det}{determinant} is $\det(AD - BC)$.

\emph{Algemene $D$.} Zij $D_t = D + tI$; dan geldt nog steeds
$CD_t = D_tC$. Beide uitdrukkingen
\[
f(t) = \det\begin{pmatrix} A & B\\ C & D_t\end{pmatrix}
\qquad\text{en}\qquad
g(t) = \det(AD_t - BC)
\]
zijn veeltermfuncties van $t$. De veelterm $\det(D + tI)$ is
monisch van graad $n$ en heeft dus hoogstens $n$ nulpunten: op
eindig veel $t$ na is $D_t$ inverteerbaar en geldt $f(t) = g(t)$
volgens het eerste geval. Twee veeltermen over een oneindig
lichaam die in oneindig veel punten overeenstemmen, zijn gelijk:
$f = g$, en $t = 0$ besluit.
\end{solution}

\begin{solution}{pb:b2:linalg:1}
\textbf{1.} $\Phi$ is lineair met $\ker\Phi = \bigcap_i
\ker\varphi_i$ (een $p$-tupel is nul dan en slechts dan als elke
component dat is). Voor de coördinaatvormen $\varepsilon_i$ van
$K^p$ geldt $\Phi^{\mathsf T}(\varepsilon_i) = \varepsilon_i \circ
\Phi = \varphi_i$, dus $\operatorname{im}\Phi^{\mathsf T} \supseteq
\operatorname{Vect}(\varphi_i)$; omgekeerd wordt
$\operatorname{im}\Phi^{\mathsf T}$ opgespannen door de
$\Phi^{\mathsf T}(\varepsilon_i)$ (de $\varepsilon_i$ spannen
$(K^p)^*$ op). Dus is $\operatorname{rk}\Phi = \operatorname{rk}
\Phi^{\mathsf T} = \dim\operatorname{Vect}(\varphi_1, \dots,
\varphi_p) =: r$ (\cref{prop:b2:linalg:transposerank}), en de
dimensiestelling geeft $\dim\bigcap_i\ker\varphi_i = n - r$.

\textbf{2.} ($\Leftarrow$) Houd een maximale vrije deelfamilie
over, zeg $\varphi_1, \dots, \varphi_r$, die dezelfde ruimte
opspant (zodat de hypothese nog steeds luidt $\bigcap_{i \leq
r}\ker\varphi_i \subseteq \ker\varphi$: de doorsnede over alle $i$
valt samen met die over $i \leq r$, want elke weggelaten vorm is
een combinatie). De afbeelding $\Psi = (\varphi_1, \dots,
\varphi_r) \colon E \to K^r$ is surjectief (vraag 1: haar rang is
$r$). Geldt $\Psi(x) = \Psi(y)$, dan is $x - y \in \ker\Psi
\subseteq \ker\varphi$, dus $\varphi(x) = \varphi(y)$: $\varphi$
factoriseert als $\varphi = \lambda \circ \Psi$ met een
welgedefinieerde $\lambda \colon K^r \to K$; en $\lambda$ is
lineair omdat $\Psi$ lineair en surjectief is (voor $t = \Psi(x)$
en $t' = \Psi(x')$: $\lambda(t + \alpha t') = \varphi(x + \alpha
x') = \lambda(t) + \alpha\lambda(t')$). Met $\lambda = \sum c_i
\varepsilon_i$ volgt $\varphi = \sum_{i \leq r} c_i\varphi_i$.
($\Rightarrow$) Is $\varphi = \sum c_i \varphi_i$, dan doodt elke
$x$ die alle $\varphi_i$ doodt, ook $\varphi$.

\textbf{3.} Volgens vraag 1 is $\dim\bigcap\ker\varphi_i = n - r$
met $r = \dim\operatorname{Vect}(\varphi_i) \leq p$, en $r = p$
dan en slechts dan als de familie vrij is. Zij $F$ een deelruimte
van codimensie $p$: haar \omterm{def:b2:linalg:annihilator}{annihilator} heeft dimensie $p$
(\cref{thm:b2:linalg:annihilator}); een basis $(\varphi_1, \dots,
\varphi_p)$ van $F^\circ$ geeft $F = \bigcap_i\ker\varphi_i$ (de
terugwinformule). Met minder gaat het niet: een doorsnede van $q$
hypervlakken heeft volgens vraag 1 dimensie $\geq n - q > n - p$.

\textbf{4.} Bereken $\ker\varphi_1 \cap \ker\varphi_2$: uit $x + y
- z = 0$ en $y + z - t = 0$ volgt, met $(y, z)$ als parameters, $x
= z - y$ en $t = y + z$, dus het vlak van de vectoren $(z - y,\;
y,\; z,\; y + z)$. Daarop is $\psi = x + 2y - t = (z - y) + 2y -
(y + z) = 0$: volgens het factorisatielemma is $\psi \in
\operatorname{Vect}(\varphi_1, \varphi_2)$ --- inderdaad is $\psi
= \varphi_1 + \varphi_2$. Maar $\psi' = x + y + t = (z - y) + y +
(y + z) = y + 2z$ is daar niet identiek nul ($y = 1, z = 0$ geeft
$1$): dus $\psi' \notin \operatorname{Vect}(\varphi_1,
\varphi_2)$.

\textbf{5.} Drie vormen op een driedimensionale ruimte: vrijheid
volstaat. Geldt $a\psi_0 + b\psi_1 + c\psi_2 = 0$, test dan op $1,
X, X^2$: $a + b + c = 0$, $b + \frac c2 = 0$, $b + \frac c3 = 0$;
de laatste twee van elkaar aftrekken geeft $c = 0$, daarna $b = 0$
en $a = 0$. Antiduale basis: met $P = \alpha + \beta X + \gamma
X^2$ en het oplossen van $\psi_i(P_j) = \delta_{ij}$ ($P(0) =
\alpha$, $P(1) = \alpha + \beta + \gamma$, $\int_0^1 P = \alpha +
\frac\beta2 + \frac\gamma3$) volgt
\[
P_0 = 1 - 4X + 3X^2, \qquad
P_1 = -2X + 3X^2, \qquad
P_2 = 6X - 6X^2 .
\]
(Controle, bijvoorbeeld: $\int_0^1 P_2 = 3 - 2 = 1$ en $P_2(0) =
P_2(1) = 0$.) Het interpolatieprobleem wordt opgelost door de
coördinaten in de antiduale basis:
\[
P = 1\cdot P_0 + 2\cdot P_1 + \tfrac32\, P_2 = 1 + X
\]
(coëfficiënt van $X$: $-4 - 4 + 9 = 1$; coëfficiënt van $X^2$: $3
+ 6 - 9 = 0$); en inderdaad is $P(0) = 1$, $P(1) = 2$, $\int_0^1 P
= \frac32$.

\textbf{6.} Lineariteit: voor elke $\varphi$ is $J(x + \alpha
y)(\varphi) = \varphi(x + \alpha y) = J(x)(\varphi) + \alpha
J(y)(\varphi)$, dat wil zeggen $J(x + \alpha y) = J(x) + \alpha
J(y)$. Injectiviteit: is $x \neq 0$, vul dan $x = e_1$ aan tot een
basis; de coördinaatvorm $e_1^*$ voldoet aan $J(x)(e_1^*) = 1 \neq
0$. Omdat $\dim E^{**} = \dim E^* = \dim E$, is injectief hier
hetzelfde als bijectief.

\textbf{7.} Insluiting: voor $x \in F$ en $\varphi \in F^\circ$ is
$J(x)(\varphi) = \varphi(x) = 0$, dus $J(F) \subseteq
F^{\circ\circ}$. Dimensies (twee keer
\cref{thm:b2:linalg:annihilator}):
\[
\dim F^{\circ\circ} = \dim E^* - \dim F^\circ
= n - (n - \dim F) = \dim F = \dim J(F),
\]
want $J$ is injectief. Bijgevolg is $J(F) = F^{\circ\circ}$.

\textbf{8.} Eerste identiteit: $\varphi$ doodt $F + G$ dan en
slechts dan als zij zowel $F$ als $G$ doodt (zij doodt sommen dan
en slechts dan als zij de delen doodt): $(F+G)^\circ = F^\circ
\cap G^\circ$. Tweede: de insluiting $F^\circ + G^\circ \subseteq
(F \cap G)^\circ$ is duidelijk (elke term doodt $F \cap G$).
Dimensies, met de eerste identiteit en Grassmann:
\[
\dim(F^\circ + G^\circ) = \dim F^\circ + \dim G^\circ -
\dim(F^\circ \cap G^\circ)
= (n - \dim F) + (n - \dim G) - \bigl(n - \dim(F +
G)\bigr),
\]
wat volgens Grassmann in $E$ gelijk is aan $n - \dim(F \cap G) =
\dim(F \cap G)^\circ$: gelijkheid.

\textbf{9.} Via het lemma: uit $\ker\psi \subseteq \ker\varphi$
met $p = 1$ volgt $\varphi \in \operatorname{Vect}(\psi)$, en
$\varphi \neq 0$ maakt de scalair ongelijk aan nul. Rechtstreeks:
kies $x_0$ met $\psi(x_0) \neq 0$; elke $x$ schrijft zich als $x =
\bigl(x - \frac{\psi(x)}{\psi(x_0)}x_0\bigr) +
\frac{\psi(x)}{\psi(x_0)} x_0$, met de eerste term in $\ker\psi =
\ker\varphi$; $\varphi$ toepassen geeft $\varphi(x) =
\frac{\varphi(x_0)}{\psi(x_0)}\psi(x)$.

\textbf{10.} Neem de \omterm{def:b2:linalg:dual}{duale basis} $(\varphi_1^*, \dots,
\varphi_n^*)$ van $(\varphi_1, \dots, \varphi_n)$ binnen $E^{**}$
(\cref{def:b2:linalg:dual} toegepast op $E^*$) en zet $u_j =
J^{-1}(\varphi_j^*)$: een basis van $E$ ($J$ is een isomorfisme,
vraag 6), met $\varphi_i(u_j) = J(u_j)(\varphi_i) =
\varphi_j^*(\varphi_i) = \delta_{ij}$. Eenduidigheid: de
voorwaarden $\varphi_i(u_j) = \delta_{ij}$ leggen $J(u_j)$ op de
basis $(\varphi_i)$ vast en dus ook $u_j$.

\textbf{11.} Lineariteit: $(u + \alpha v)^{\mathsf T}\psi = \psi
\circ (u + \alpha v) = u^{\mathsf T}\psi + \alpha\, v^{\mathsf
T}\psi$. Injectiviteit: is $u \neq 0$, kies dan $x$ met $u(x) \neq
0$ en $\psi$ met $\psi(u(x)) \neq 0$ (de truc met de
coördinaatvorm uit vraag 6): dan is $u^{\mathsf T}\psi \neq 0$. De
ruimten $\mathcal{L}(E,F)$ en $\mathcal{L}(F^*, E^*)$ hebben beide
dimensie $\dim E \dim F$: dus bijectief. Is $u$ inverteerbaar, dan
geeft de omkeringsregel $(vu)^{\mathsf T} = u^{\mathsf
T}v^{\mathsf T}$ dat $u^{\mathsf T}(u^{-1})^{\mathsf T} =
(u^{-1}u)^{\mathsf T} = \mathrm{id}_{E^*}$ en
$(u^{-1})^{\mathsf T}u^{\mathsf T} = (uu^{-1})^{\mathsf T} =
\mathrm{id}_{F^*}$, dus $(u^{\mathsf T})^{-1} = (u^{-1})^{\mathsf
T}$.

\textbf{12.} Voor $x \in E$ en $\psi \in F^*$:
\[
\bigl(u^{\mathsf T\mathsf T}(J_E x)\bigr)(\psi)
= (J_E x)\bigl(u^{\mathsf T}\psi\bigr)
= (u^{\mathsf T}\psi)(x)
= \psi\bigl(u(x)\bigr)
= \bigl(J_F(u(x))\bigr)(\psi).
\]
Omdat $\psi$ willekeurig is, geldt $u^{\mathsf T\mathsf T} \circ
J_E = J_F \circ u$.

\textbf{13.} Volgens \cref{prop:b2:linalg:transposerank} is $\ker
u^{\mathsf T} = (\operatorname{im} u)^\circ$, dus is $u$
surjectief $\iff \operatorname{im} u = F \iff
(\operatorname{im}u)^\circ = \{0\}$
(\cref{thm:b2:linalg:annihilator}) $\iff u^{\mathsf T}$ injectief.
En $\operatorname{im} u^{\mathsf T} = (\ker u)^\circ$, dus is $u$
injectief $\iff \ker u = \{0\} \iff (\ker u)^\circ = E^*$ $\iff
u^{\mathsf T}$ surjectief.

\textbf{14.} Is $u(F) \subseteq F$ en $\varphi \in F^\circ$, dan is
$(u^{\mathsf T}\varphi)(x) = \varphi(u(x)) = 0$ voor $x \in F$,
dus $u^{\mathsf T}\varphi \in F^\circ$. Omgekeerd, is $u(F)
\not\subseteq F$, kies dan $x \in F$ met $u(x) \notin F$; volgens
de terugwinformule van \cref{thm:b2:linalg:annihilator} is er een
$\varphi \in F^\circ$ met $\varphi(u(x)) \neq 0$: dan is
$(u^{\mathsf T}\varphi)(x) \neq 0$ hoewel $x \in F$, zodat
$u^{\mathsf T}\varphi \notin F^\circ$ en $F^\circ$ niet stabiel
is.

\textbf{15.} Er geldt $u^{\mathsf T} - \lambda\,\mathrm{id}_{E^*}
= (u - \lambda\,\mathrm{id}_E)^{\mathsf T}$ (transponeren is
lineair en $\mathrm{id}^{\mathsf T} = \mathrm{id}$), dus is de kern
gelijk aan $(\operatorname{im}(u - \lambda\,\mathrm{id}))^\circ$
(\cref{prop:b2:linalg:transposerank}), van dimensie
\[
n - \operatorname{rk}(u - \lambda\,\mathrm{id})
= \dim\ker(u - \lambda\,\mathrm{id})
\]
volgens de dimensiestelling. In het bijzonder is de ene kern
ongelijk aan nul precies wanneer de andere dat is: dezelfde
eigenwaarden, dezelfde meetkundige multipliciteiten.

\textbf{16.} Insluiting: is $b = u(x)$ en $u^{\mathsf T}\psi = 0$,
dan is $\psi(b) = \psi(u(x)) = (u^{\mathsf T}\psi)(x) = 0$; dus
$\operatorname{im} u \subseteq (\ker u^{\mathsf T})_\circ$.
Dimensies: voor een deelruimte $S \subseteq F^*$ is $S_\circ =
J_F^{-1}(S^\circ)$ (werk het uit: $y \in S_\circ$ dan en slechts
dan als elke $\psi \in S$ het element $y$ doodt, dan en slechts
dan als $J_F(y) \in S^\circ$), dus $\dim S_\circ = \dim F - \dim
S$. Met $S = \ker u^{\mathsf T}$:
\[
\dim(\ker u^{\mathsf T})_\circ
= \dim F - \dim\ker u^{\mathsf T}
= \operatorname{rk} u^{\mathsf T} = \operatorname{rk} u :
\]
de dimensies zijn gelijk, dus $\operatorname{im} u = (\ker
u^{\mathsf T})_\circ$. Anders gezegd: $b \in \operatorname{im} u$
dan en slechts dan als $\psi(b) = 0$ voor elke $\psi$ met
$u^{\mathsf T}\psi = 0$ --- het alternatief van Fredholm.

\textbf{17.} Vereenzelvig $(K^m)^*$ met $K^m$ via $y \mapsto
\psi_y$, $\psi_y(v) = y^{\mathsf T}v$; dan is $(u^{\mathsf
T}\psi_y)(x) = y^{\mathsf T}Ax = (A^{\mathsf T}y)^{\mathsf T}x$,
dus $u^{\mathsf T}\psi_y = \psi_{A^{\mathsf T}y}$: de
\omterm{def:b2:linalg:transpose}{getransponeerde afbeelding} is de \omterm{def:b2:linalg:transpose}{getransponeerde} matrix.
\emph{Hoogstens één:} geldt $Ax = b$ en $A^{\mathsf T}y = 0$, dan
is $y^{\mathsf T}b = y^{\mathsf T}Ax = (A^{\mathsf T}y)^{\mathsf
T}x = 0 \neq 1$. \emph{Minstens één:} faalt (i), dan levert vraag
16 een $\psi_y$ met $A^{\mathsf T}y = 0$ en $y^{\mathsf T}b \neq
0$; schaal $y$ zo dat dit $1$ wordt.

\textbf{18.} Hier is $A = \begin{psmallmatrix} 1 & 1 & 0\\ 0 & 1 &
1\\ 1 & 2 & 1\end{psmallmatrix}$ (derde rij = eerste + tweede, dus
$A$ is singulier). Los $A^{\mathsf T}y = 0$ op: $y_1 + y_3 = 0$,
$y_1 + y_2 + 2y_3 = 0$, $y_2 + y_3 = 0$ geven $y_1 = y_2 = -y_3$:
de rechte opgespannen door $y = (1, 1, -1)$. Fredholm: oplosbaar
dan en slechts dan als $y^{\mathsf T}b = b_1 + b_2 - b_3 = 0$, dat
wil zeggen $b_3 = b_1 + b_2$ --- zichtbaar de juiste voorwaarde,
want de derde vergelijking is de som van de eerste twee.

\textbf{19.} De matrix van $L$ heeft $1$ op de diagonaal en
$-\frac12$ op de plaatsen $(k, k\pm1)$ (modulo $n$): symmetrisch,
dus $L^{\mathsf T} = L$ onder de vereenzelviging van vraag 17.
\emph{Kern:} is $Lx = 0$, dan geldt $x_k = \frac12(x_{k-1} +
x_{k+1})$ voor elke $k$. Zij $k_0$ een index waar $x_k$ maximaal
is; het gemiddelde van de twee buren, die beide $\leq x_{k_0}$
zijn, is alleen gelijk aan $x_{k_0}$ wanneer beide gelijk zijn aan
$x_{k_0}$; door de \omterm{def:b2:structures:sn}{cykel} rond voort te planten volgt dat $x$
constant is. Omgekeerd worden constanten gedood. Dus $\ker
L^{\mathsf T} = \ker L = \R(1, \dots, 1)$, en het alternatief van
Fredholm luidt: $Lx = b$ is oplosbaar dan en slechts dan als
$(1,\dots,1)^{\mathsf T} b = \sum_k b_k = 0$ --- de discrete
verenigbaarheidsvoorwaarde: een ``warmteverdeling'' op een ring is
door een potentiaal te realiseren precies wanneer haar totale flux
nul is.

\textbf{20.} Is $A$ antisymmetrisch en $S$ symmetrisch, dan
\[
\operatorname{tr}(AS) = \operatorname{tr}\bigl((AS)^{\mathsf
T}\bigr) = \operatorname{tr}(S^{\mathsf T}A^{\mathsf T}) =
-\operatorname{tr}(SA) = -\operatorname{tr}(AS),
\]
dus $2\operatorname{tr}(AS) = 0$ en (want $\operatorname{char} K
\neq 2$) $\operatorname{tr}(AS) = 0$: $\mathcal{A}_n \subseteq
\mathcal{S}_n^\circ$ (met de vereenzelviging van de \omterm{def:b2:linalg:dual}{duale ruimte}
met matrices). Dimensies: $\dim\mathcal{S}_n^\circ = n^2 -
\frac{n(n+1)}2 = \frac{n(n-1)}2 = \dim\mathcal{A}_n$: gelijkheid.
Met verwisselde rollen (dezelfde berekening) volgt
$\mathcal{A}_n^\circ = \mathcal{S}_n$.

\textbf{21.} Er geldt $\operatorname{tr}(I_nM) = \operatorname{tr}
M = 0$ voor $M \in \mathfrak{sl}_n$: de rechte $KI_n$ ligt in de
\omterm{def:b2:linalg:annihilator}{annihilator}, waarvan de dimensie $n^2 - (n^2 - 1) = 1$ is:
gelijkheid. Vertaald door het isomorfisme $A \mapsto
\operatorname{tr}(A\,\cdot)$: een vorm die op $\mathfrak{sl}_n$
verdwijnt, is $\operatorname{tr}(\lambda I_n\,\cdot) =
\lambda\operatorname{tr}$.

\textbf{22.} Zij $M \in \mathcal{M}_n(K)$. De veelterm $t \mapsto
\det(M - tI)$ is ongelijk aan nul en van graad $n$, dus heeft zij
hoogstens $n$ nulpunten; $K$ heeft karakteristiek $0$ en is dus
oneindig: kies een $\lambda \neq 0$ die geen nulpunt is. Dan
schrijft $M = (M - \lambda I) + \lambda I$ de matrix $M$ als som
van twee inverteerbare matrices.

\textbf{23.} \emph{Stap 1:} pas voor inverteerbare $P$ en
willekeurige $X$ de invariantie toe op $M = XP$: $t(P(XP)P^{-1}) =
t(XP)$, dat wil zeggen $t(PX) = t(XP)$. \emph{Stap 2:} leg $X$
vast; beide leden van $t(BX) = t(XB)$ zijn lineair in $B$ en
stemmen overeen op inverteerbare $B$; volgens vraag 22 is elke $B$
een som van twee inverteerbare matrices, dus stemmen zij overal
overeen. \emph{Stap 3:} $t$ doodt elke commutator $XB - BX$; de
commutatoren spannen $\mathfrak{sl}_n$ op (aangetoond in het
bewijs van \cref{prop:b2:linalg:trace}), dus verdwijnt $t$ op
$\mathfrak{sl}_n$ en geeft vraag 21 dat $t = c\operatorname{tr}$.
(Omgekeerd is elke $c\operatorname{tr}$ invariant onder
gelijkvormigheid: het spoor is \emph{de} lineaire invariant van de
gelijkvormigheid.)

\textbf{24.} Is $\operatorname{rk} u = r' \leq r$, neem dan een
basis $(f_1, \dots, f_{r'})$ van $\operatorname{im} u$ en schrijf
$u(x) = \sum_{i=1}^{r'} \psi_i(x) f_i$; elke coördinaat
$\psi_i(x)$ van $u(x)$ is lineair in $x$ (de samenstelling van $u$
met een coördinaatvorm), dus is $u$ een som van $r' \leq r$
afbeeldingen van rang $\leq 1$ (vul aan met nullen). Omgekeerd, is
$u = \sum_{i=1}^r \psi_i(\cdot)f_i$, dan is $\operatorname{im} u
\subseteq \operatorname{Vect}(f_1, \dots, f_r)$, dus
$\operatorname{rk} u \leq r$. Subadditiviteit: schrijf $u$ met
$\operatorname{rk} u$ termen en $v$ met $\operatorname{rk} v$
termen; de som heeft $\operatorname{rk} u + \operatorname{rk} v$
termen, dus $\operatorname{rk}(u + v) \leq \operatorname{rk} u +
\operatorname{rk} v$.

\textbf{25.} Het woordenboek: een deelruimte $F$ correspondeert met
$F^\circ$ met complementaire dimensie
(\cref{thm:b2:linalg:annihilator}), en terug via de bidualiteit
(vragen 6--7); sommen wisselen met doorsneden (vraag 8); een
afbeelding $u$ correspondeert met $u^{\mathsf T}$, met $\ker
u^{\mathsf T} = (\operatorname{im}u)^\circ$,
$\operatorname{im}u^{\mathsf T} = (\ker u)^\circ$, gelijke rangen,
verwisselde injectiviteit en surjectiviteit, en overeenkomende
stabiele deelruimten en eigenwaarden (vragen 11--15); de
vergelijking $u(x) = b$ is oplosbaar dan en slechts dan als $b$
loodrecht staat op $\ker u^{\mathsf T}$ (vragen 16--19); en op
$\mathcal{M}_n$ maakt de spoorparing het hele woordenboek
concreet, met het spoor als de enige lineaire invariant van de
gelijkvormigheid (vragen 20--23) en de rang als de minimale lengte
van een ontbinding in elementaire tensoren (vraag 24). In
oneindige dimensie falen de dimensietellingen en worden zij
vervangen door hypothesen over de \emph{geslotenheid} van beelden
en door volledigheid --- op hilbertruimten wordt dit de
representatiestelling van Riesz en de theorie van Fredholm voor
compacte operatoren, in het volume van bachelorjaar 3 eerlijk
bewezen.
\end{solution}
